\section{Actual slow-band transport in the two character families}
\label{sec:character-slow-band}

\paragraph{MA-SB-004 (the character slow-band limit with its exact remainder).}
Let $\eta$ be a nonprincipal primitive Dirichlet character, or
the non-$\mathbb Q$-induced cubic character $\psi_{269,3}$ of
MA-CH-003.  More generally the statement applies to the
nontrivial finite-order characters of $\mathbb Q(\zeta_3)$
covered by MA-CGP-001--004.  Use probability Haar measure
on the norm-one compact group, the coherent additive
character of MA-NR-001, and the primitive finite factors
with $c_v=\beta(p_v)^{-1}$ as in MA-CH.3.
Fix
\[
 \begin{gathered}
 0<\alpha<\tfrac12,\qquad s=0,\qquad
 g\in C_c^\infty((0,\infty)),\qquad h(t)=g(\log t),\\
 G_0(z)=\int g(x)e^{zx}\,dx,\qquad
 Q=|D_K|N\mathfrak f_\eta,\qquad Y=Q\Lambda\ge1.
\end{gathered}
 \tag{MA-SB.30}
\]
Thus $h(1)=0$.  Let $Q_{\eta,Y}^{\rm glob}$ be the actual
orthogonal PW projection at invariant source product level
$Y$, with $S'=S$ and $S$ containing the conductor and
different.  Equivalently it is the primitive archimedean
norm-band projection of MA-CGP-002 at $\Lambda=Y/Q$.
No Euler factor has been removed by this equality.

Let $\rho_{X,\delta}$ be a source cutoff supported on
$[1/X,X]$, equal to one on
$[1/X+\delta,X-\delta]$, and between zero and one.
There are two fixed functions
$\omega_{\eta,g},\nu_{\eta,g}:[0,\infty)\to[0,\infty)$,
both tending to zero at infinity, such that whenever
$X\ge2Y\ge2$ and $0<\delta\le1/(4X)$,
\[
 \begin{split}
 &\left|\operatorname{Tr}_H
      (M_{\rho_{X,\delta}}C_gQ_{\eta,Y}^{\rm glob})
                +\tau_{\eta,+}(g)\right|\\
 &\hspace{15mm}\le
    \omega_{\eta,g}(\log(X/2))
        +\nu_{\eta,g}(\log(X/2)),\\
 &\tau_{\eta,+}(g)
     =\sum_{\operatorname{Re}\rho_\eta<1/2+\alpha}
                 m_\eta(\rho_\eta)G_0(\rho_\eta-1/2).
\end{split}
 \tag{MA-SB.31}
\]
The trace is on the actual character-sector space
$H=L^2(\mathbb R,e^{2\alpha|x|}\,dx)$.
In particular its limit along every source-thin path
$X,Y\to\infty$, $Y/X\to0$, is $-\tau_{\eta,+}(g)$.

For a comparison with the actual local quotient, use the
common fixed profiles from MA-CQ-003:
\[
 \begin{aligned}
 a_X(t)&=A(t/X)-A(Xt),\\
 T^{\beta,S}_{X,Y,X^{-1}}(h)
  &=M_{a_X}A_hD_{\beta,S}Q_{\infty,Y}D_{\beta,S}^*,\\
 Q_{\infty,Y}
  &=\mathcal F_\infty^{-1}
       M_{B(|\xi|_\infty/Y)}\mathcal F_\infty.
\end{aligned}
 \tag{MA-SB.32}
\]
Here $A,B$ equal one on $[0,1]$ and zero on $[2,\infty)$,
and $M_{a_X}$ acts through the archimedean norm in the
local operator and through $t=e^x$ in the global operator.
The finite unitary is the actual
$D_{\beta,S}=M_{\bar\beta,S}$ of MA-CQ.16--.17;
the intermediate Fourier map exchanges the dual sectors.
Choose the fixed $S$ also to include every prime ideal
whose positive or inverse prime-power norm meets
$\operatorname{supp}h$.
Then the two actual limits exist and their difference is
\[
 \begin{split}
 &\lim_{\substack{X,Y\to\infty\\Y/X\to0}}
 \left\{\operatorname{Tr}_H(M_{a_X}C_gQ_{\eta,Y}^{\rm glob})
            -\operatorname{Tr}T^{\beta,S}_{X,Y,X^{-1}}(h)\right\}\\
 &\hspace{12mm}=
    \sum_{\operatorname{Re}\rho_\eta\ge1/2+\alpha}
                  m_\eta(\rho_\eta)G_0(\rho_\eta-1/2).
\end{split}
 \tag{MA-SB.33}
\]
There is no polar term in these nontrivial primitive
character sectors.  The upper boundary is included in
the remainder.  Neither its nonemptiness nor its value
on an arbitrary particular test is asserted.

\begin{proof}
\emph{The actual range and its normalization.}
By MA-CGP-002 the source range is the actual centered
character ideal supported on
$x\ge-\log(Q\Lambda)=-\log Y$.
Put $L=\log Y$ and $r=x+L$.
The unitary $f\mapsto e^{\alpha x}f(x)$, followed by
this translation, gives exactly the metric
\[
 R_L(r)=
 \begin{cases}e^{4\alpha(L-r)},&0\le r\le L,\\
 1,&r\ge L.
 \end{cases}
 \tag{MA-SB.34}
\]
In the equivalent ordinary half-line norm the range is
$V_\eta L^2(0,\infty)$, where
\[
 V_\eta=\mathcal L^{-1}M_{B_{\eta,\alpha}}\mathcal L,
 \qquad P_{K_\eta}=1-V_\eta V_\eta^*.
 \tag{MA-SB.35}
\]
The actual pure Blaschke divisor consists of
$z=\alpha+1/2-\rho_\eta$ with positive real part,
including the full multiplicities.  This is the
$\eta$ divisor: the $\bar\eta$ in the reflected
Fourier formula MA-CGP.13 has already been changed
by the completed functional equation.  Neither
complex conjugation of the final divisor nor a
further conductor dilation is allowed here.

Put $k(t)=e^{\alpha t}g(t)$.
MA-CGP-004, at $c=0$ where the metric is one, proves
the actual ambient lift
$C_k^+P_{K_\eta}\in\mathcal S_1$ and identifies its
trace as $\tau_{\eta,+}(g)$.
This verifies the operator hypothesis of MA-SB-001;
a summable formal spectrum would not suffice.
The proof there supplies the fixed trace-class
Hankel operator $H_{B_\eta k}$ and hence the tails
\[
 \begin{aligned}
 \omega_{\eta,g}(r)&=
       \|P_{[r,\infty)}C_k^+P_{K_\eta}\|_1,\\
 \nu_{\eta,g}(r)&=
       \|P_{[r,\infty)}H_{B_\eta k}\|_1.
\end{aligned}
 \tag{MA-SB.36}
\]
Both vanish at infinity with no separation of the
zeros from the half-plane boundary.

\emph{The actual weighted Gram trace.}
The precise weighted projector is
\[
 V_\eta\mathfrak G_{\eta,L}^{-1}V_\eta^*M_{R_L},
 \qquad
 \mathfrak G_{\eta,L}=V_\eta^*M_{R_L}V_\eta .
 \tag{MA-SB.37}
\]
Causality and $R_L-1$ supported on $[0,L]$ give
$D_{\eta,L}=\mathfrak G_{\eta,L}^{-1}-1
 =P_{[0,L]}D_{\eta,L}P_{[0,L]}$ and
$-1\le D_{\eta,L}\le0$.
Every hypothesis used in MA-SB.18--.22 is now
verified for this fixed $V_\eta$ and $k$.
In particular those identities retain the right
metric factor, establish trace class before
cyclicity, and insert the input support projection
only inside an already defined trace-class trace.
They prove the same error bound
$\omega_{\eta,g}(R)+\nu_{\eta,g}(R-L)$.

For the source cutoff, $X\ge2Y$ and the stated
$\delta$ remove its entire lower transition from
the support half-line.  Its half-line profile is
one up to $R=L+\log(X-\delta)$.
This proves (MA-SB.31), including uniformity in
$1\le Y\le X/2$.
For $a_X$, the lower summand vanishes on that
half-line and the upper profile is one through
$R=L+\log X$.  Thus the same actual global limit
holds for the profile in (MA-SB.32).

\emph{The actual local limit and its trace comparison.}
MA-CQ-003 applies with $\epsilon=X^{-1}$:
$XY\to\infty$ and $\epsilon Y=Y/X\to0$.
It gives, on its actual local quotient completion,
\[
 \operatorname{Tr}T^{\beta,S}_{X,Y,X^{-1}}(h)
       \longrightarrow A_\infty(h)+B_{\beta,S}(h).
 \tag{MA-SB.38}
\]
The parity or complex archimedean coefficient and
the primitive dual finite factors are precisely
those of MA-CQ.21.  The conductor/different
dilations in $D_{\beta,S}$ remain in this operator;
their scalar trace contribution is proportional
to $g(0)=0$, as proved in MA-CQ.23.
The choice of $S$ gives $B_{\beta,S}=B_\eta$.

Meyer's actual smooth quotient trace and its
character-sector comparison MA-CC-001 identify
$A_\infty+B_\eta$ with
$-\tau_\eta^{\rm full}(g)$.  The polar quotient
is zero because the character is nontrivial.
This uses the unconditional smooth comparison,
not its separate uniform-interior Hilbert version.
Subtracting the two already proved operator
limits gives $\tau_\eta^{\rm full}-\tau_{\eta,+}$;
their established divisor descriptions then give
(MA-SB.33), with its upper-edge convention.
\end{proof}

The comparison with the actual completed full-strip
trace of MA-WP is also exact:
\[
 \begin{split}
 &\lim_{\substack{X,Y\to\infty\\Y/X\to0}}
    \operatorname{Tr}(M_{a_X}C_gQ_{\eta,Y}^{\rm glob})
       -\left(-\tau_{\eta,\alpha}(g)\right)\\
 &\hspace{10mm}=-\sum_{\operatorname{Re}\rho_\eta\le1/2-\alpha}
                  m_\eta(\rho_\eta)G_0(\rho_\eta-1/2).
 \end{split}
 \tag{MA-SB.39}
\]
At $s=0$ no boundary jet is continuous, which explains
the included lower edge on the right.
Every finite global trace here has the concrete
character Gram-resolvent realization MA-CGP.18.
At fixed $X,Y$, take that proved regularizer limit
first, or choose a regularizer whose fixed-data
trace error tends to zero along the path.
No uniform regularizer rate, LF lift of its inverse,
or omission of either Gram factor is asserted.

For $\psi_{269,3}$ the scale in (MA-SB.30) is
$Q=3\cdot269^2$, and its actual local finite
unitary is the one in MA-CQ-004.  In particular the
prime over $3$ retains its Euler factor, whereas
the ramified prime over $269$ contributes its
conductor dilation.  This distinction is already
part of the proved operator, not a modification
of the primitive $L$-function.
The result uses neither $\mathsf A$ nor $M$.
Its scope remains $s=0$, the stated positive
logarithmic support, and the actual source level
$S'=S$.  The paired local comparison uses the
specified fixed profiles, without an unproved
uniform claim for arbitrary varying local smoothings.
